\chapter{Análisis del sistema de información}

\section{Contexto y usuarios}
\pending{Identificar investigadores, estudiantes, desarrolladores, tutores y otros perfiles relevantes.}

\section{Problema del flujo experimental actual}
\pending{Describir la dispersión de scripts, configuraciones y resultados que motiva la plataforma.}

\section{Ciclo de vida de un experimento}
\pending{Analizar creación, configuración, validación, ejecución, seguimiento, evaluación, persistencia y comparación.}

\section{Casos de uso}
\pending{Especificar casos de uso para gestionar modelos, datasets, experimentos, ejecuciones, resultados y exportaciones.}
\placeholderfigure{Diagrama de casos de uso}{6cm}

\section{Requisitos funcionales}
\pending{Definir requisitos verificables de gestión de usuarios, modelos, datasets, ejecuciones, resultados y comparación.}

\section{Requisitos no funcionales}
\pending{Definir requisitos de reproducibilidad, modularidad, extensibilidad, usabilidad, rendimiento, trazabilidad y mantenibilidad.}

\section{Datos gestionados}
\pending{Identificar datasets, configuraciones, modelos, ejecuciones, métricas, logs y artefactos.}

\section{Restricciones}
\pending{Documentar simulación local, CPU, parameter-shift, shots, ruido, número de qubits y alcance inicial.}

\section{Alcance funcional}
\pending{Separar capacidades generales de la plataforma, capacidades específicas de la QCNN y extensiones futuras.}

\section{Criterios de aceptación}
\pending{Relacionar cada requisito con una comprobación, una prueba o un resultado observable.}
